{\LARGE\bfseries Six Elements, One Comparison}\par
{\large A teaching companion to \emph{Quotients, Refinement and Reconstruction}}\par
{\small October 2026}

This booklet starts with addition modulo six and follows one example through
the comparison argument behind Jordan--Hölder. You need subgroups and cosets;
the categorical language is introduced only after its job becomes visible.
Read §§1--8 as the main route. Section 9 is an optional spectral calculation,
§§10--11 contain exercises and answers, and §12 is a guide to the higher chapters.
References marked \emph{main §} lead to the full exposition.

\section{Two chains through the same six elements}\label{c:chains}
Write \(C_d=\Z/d\Z\). Let
\[
 G=\Z/6\Z,\qquad A=\{0,3\},\qquad B=\{0,2,4\}.
\]
All arithmetic in \(G\) is modulo six. We write \(0\) for the subgroup
\(\{0\}\), which has one element. Consider
\[
\begin{tikzcd}[column sep=large,row sep=small]
0\arrow[r,hook]&A\arrow[r,hook]&G\\
0\arrow[r,hook]&B\arrow[r,hook]&G.
\end{tikzcd}
\]
Every subgroup here is normal because \(G\) is abelian. The factors are
\[
\begin{array}{c|cc}
 &\text{first step}&\text{second step}\\ \hline
 A\text{-chain}&A/0\cong C_2&G/A\cong C_3\\
 B\text{-chain}&B/0\cong C_3&G/B\cong C_2 .
\end{array}
\]
These identifications are explicit: send \(3\in A\) to \(1\in C_2\),
send \(2\in B\) to \(1\in C_3\), and use reduction modulo three and modulo
two for the two quotients of \(G\). Every nonzero element of \(C_2\) or
\(C_3\) generates its group, so these factors are simple.

We can already see the same factors in a different order. Our task is to
construct the comparison that makes this inevitable, without assuming that
the number of factors is invariant.

Two elementary observations will do much of the work:
\[
 A\cap B=0,\qquad A+B=G.
\]
The first follows from the displayed lists. For the second, adding \(0\) to
the elements of \(B\) gives \(0,2,4\), and adding \(3\) gives \(3,5,1\).
Thus every element has an expression \(a+b\). This expression is unique:
if \(a+b=a'+b'\), then \(a-a'=b'-b\in A\cap B=0\).

\pausehere{Neither \(A\subseteq B\) nor \(B\subseteq A\). A single subgroup
chain cannot contain both. We will construct two refinements and match their
factors.}
\route{Main route: \mainref{3}{main §3}, the general refinement theorem,
and \mainref{4-4}{main §4.4}, the \(C_6\) comparison.}

\clearpage
\section{An algebraic maneuver is a map}\label{c:maps}
Start with \(a=b+c\). To recover \(b\), translate by \(-c\).
The translation \(\tau_c(x)=x+c\) is a bijection of the underlying set, with
inverse \(\tau_{-c}\), because associativity gives
\[
 (x+c)+(-c)=x+(c-c)=x.
\]
Thus subtraction is the application of an inverse map. A translation is
usually not a group homomorphism: it need not preserve zero.

Distributivity says something different. The homomorphism \(h(x)=2x\)
satisfies
\[
 h(x+c)=h(x)+h(c),\qquad
 \begin{tikzcd}[column sep=large,row sep=small]
 G\arrow[r,"\tau_c"]\arrow[d,"h"']&G\arrow[d,"h"]\\
 G\arrow[r,"\tau_{h(c)}"']&G .
 \end{tikzcd}
\]
The square says \(h\circ\tau_c=\tau_{h(c)}\circ h\): applying \(h\) before
or after the appropriate translation gives the same result.
For a vector space the same statement holds for multiplication by a scalar;
for a nonzero scalar that linear map is also invertible.

Now consider the quotient map \(q_A:G\to G/A\), \(x\mapsto x+A\).
Its fibers are
\[
 \{0,3\},\qquad\{1,4\},\qquad\{2,5\}.
\]
Addition of cosets is well defined: replacing \(x,y\) by \(x+a,y+a'\),
with \(a,a'\in A\), changes \(x+y\) by the element \(a+a'\) of \(A\).

Suppose \(h:G\to H\) is a group homomorphism constant on these fibers.
Then there is a unique homomorphism \(\bar h:G/A\to H\) such that
\[
 \bar h(x+A)=h(x),\qquad
 \begin{tikzcd}[column sep=large,row sep=small]
 G\arrow[r,"h"]\arrow[d,"q_A"']&H\\
 G/A\arrow[ur,dashed,"\bar h"']&
 \end{tikzcd}
\]
The fiber condition proves independence of the representative; the
homomorphism law follows from that of \(h\); surjectivity of \(q_A\) proves
uniqueness. This is \emph{descent through a quotient}.

For a concrete instance, \(h(x)=4x\) lands in \(B\) and kills \(3\).
Consequently
\[
 \beta:B\longrightarrow G/A,\quad b\longmapsto b+A,
 \qquad
 \gamma:G/A\longrightarrow B,\quad x+A\longmapsto4x
\]
are inverse homomorphisms. Indeed, \(4b=b\) for \(b\in B\), and
\(4x-x=3x\in A\) for every \(x\in G\).

The structural viewpoint records which map is invertible, which equation
expresses compatibility, and which quotient allows a map to descend.
Associativity and distributivity supply the checks; they have not disappeared.
\pausehere{Why does \(x+A\mapsto x\) fail to define a map \(G/A\to G\)?
Compare the representatives \(0\) and \(3\).}
\route{Further development: \mainref{1}{main §1}, relations, images and quotients.}

\clearpage
\section{Keep the identifications, including where they live}\label{c:relations}
A relation \(R\) from \(X\) to \(Y\) is a set of pairs \((x,y)\).
Write \(xRy\) for membership. Its \emph{input support} consists of the
\(x\)'s occurring in a pair; its \emph{output support} consists of the
\(y\)'s occurring in a pair. A relation may have several outputs or none.

For a subgroup \(D\leq G\), write
\(\Delta_D=\{(d,d):d\in D\}\). Encode our four quotient steps by
\[
\begin{array}{c|l|c|c}
 \text{relation}&\text{rule}&\text{support}&\text{class of }0\\ \hline
 r_1&x=y\in A&A&0\\
 r_2&y-x\in A&G&A\\
 s_1&x=y\in B&B&0\\
 s_2&y-x\in B&G&B .
\end{array}
\]
Thus \(r_1=\Delta_A\) records \(A\to A/0\), whereas \(r_2\) records
\(G\to G/A\): two elements are related exactly when their quotient images
agree. Each is an equivalence relation on its support. Viewed on all of
\(G\), it is a \emph{partial equivalence relation}. For instance,
\(1\) is not related to itself by \(r_1\).

Let \(Q(r)\) mean the support modulo this equivalence relation. Then
\(Q(r_1)=A\) and \(Q(r_2)=G/A\). In groups, a congruence is determined
by its identity class. The more general argument retains the whole relation,
so it need not assume this group-specific fact.

Our convention is that \(SR\) means \emph{first \(R\), then \(S\)}:
\[
 x(SR)z\quad\Longleftrightarrow\quad
 \text{there is }y\text{ with }xRy\text{ and }ySz.
\]
The opposite relation \(R^\op\) reverses every pair. Associativity of
composition is now a statement about witnesses: both \(T(SR)\) and
\((TS)R\) relate \(x\) to \(w\) precisely when there are \(y,z\) with
\(xRy,\ ySz,\ zTw\). Regrouping changes the parentheses, not the data.

For the group argument we use \emph{internal relations}: subgroups
\(R\leq X\times Y\), not arbitrary subsets. Such a relation has the
three-corner property
\[
 xRy,\quad x'Ry,\quad x'Ry'
 \quad\Longrightarrow\quad xRy'.
\]
Indeed, closure under multiplication and inversion gives
\[
 (x,y)(x',y)^{-1}(x',y')=(x,y').
\]
The order displayed works even in noncommutative groups. This property is
called \emph{difunctionality}; it is the key relation identity supplied by
the Mal'cev framework.

\pausehere{The three-corner property does not require a unique output.
It says that three corners of a related rectangle force the fourth.}
\route{Main §§1--2 develop the same operations using pullbacks, regular
images and effective quotients, so that they apply beyond groups.}

\clearpage
\section{A comparison with two outputs per input}\label{c:one-cell}
Before filling a whole comparison table, examine the relation
\[
 R=r_2s_1,\qquad bRy\quad\Longleftrightarrow\quad
 b\in B\ \text{and}\ y-b\in A.
\]
First \(s_1\) requires the input to be an element of \(B\) and leaves it
fixed. Then \(r_2\) allows movement by an element of \(A\). Its six pairs are
displayed by
\[
\begin{array}{c|c}
 \text{input}&\text{possible outputs}\\ \hline
 0&0,\ 3\\
 2&2,\ 5\\
 4&4,\ 1 .
\end{array}
\]
The input support is \(B\); the output support is all of \(G\).
There is no single output in the definition.

There are two kinds of ambiguity to remove. Define
\[
 E=R^\op R,\qquad F=RR^\op.
\]
Two inputs are \(E\)-related when they share an output. The three rows of
the table are disjoint, so \(E=\Delta_B\). Two outputs are \(F\)-related
when they share an input. Hence \(F=r_2\), with classes
\(\{0,3\},\{2,5\},\{4,1\}\).

After quotienting, the two possible outputs in each row become one:
\[
 B/E=B\xrightarrow[\cong]{\ \beta\ }G/F=G/A,\qquad b\longmapsto b+A.
\]
We proved the inverse formula \(x+A\mapsto4x\) in §\ref{c:maps}.
The relation therefore produces an actual isomorphism, with no choice
of representatives in its definition.

Here is the same construction as a diagram. Let \(M=R\), regarded as a
subgroup of \(B\times G\), and take its coordinate projections
\(a(b,y)=b\) and \(b'(b,y)=y\):
\[
\begin{tikzcd}[row sep=large,column sep=large]
 &M\arrow[dl,two heads,"a"']\arrow[dr,two heads,"b'"]&\\
 B\arrow[d,"\id_B"']&&G\arrow[d,"q_A"]\\
 B\arrow[rr,"\beta"']&&G/A .
\end{tikzcd}
\]
Both projections are onto their supports. The equation
\(\beta a=q_A b'\) says that related elements receive matching quotient
images. It characterizes \(\beta\) uniquely because \(a\) is surjective.

\pausehere{The pair \((2,5)\) belongs to \(R\). Does \((5,2)\)?
Keep input support and output support separate.}
\route{This is the concrete instance of the common-quotient construction
in \mainref{2}{main §2}.}

\clearpage
\section{Why the construction works beyond six elements}\label{c:general}
Let \(R\leq X\times Y\) be any subgroup relation. Restrict to its input
and output supports \(X'\leq X\), \(Y'\leq Y\). Set
\(E=R^\op R\) and \(F=RR^\op\).
They are reflexive on these supports and symmetric by definition.

For transitivity of \(E\), suppose \(x,x'\) share output \(y\), and
\(x',x''\) share output \(z\). The three corners
\(xRy,\ x'Ry,\ x'Rz\) give \(xRz\), so \(x,x''\) share \(z\).
Interchanging input and output proves transitivity of \(F\).
Thus \(E,F\) are equivalence relations on their supports.

They respect the group operations. For example, witnesses \(y,v\) for
\(xEx'\) and \(uEu'\) give witness \(yv\) for
\((xu)E(x'u')\); inversion of a witness gives the corresponding inverse
relation. The same holds for \(F\). Their classes therefore carry
well-defined quotient group structures.

Define
\[
 \theta:X'/E\longrightarrow Y'/F,\qquad [x]_E\longmapsto[y]_F
 \quad\text{whenever }xRy.
\]
An output exists by the definition of support. Different outputs for the
same \(x\) are \(F\)-related. If \(xEx'\), choose a common output \(y_0\);
any outputs for \(x\) and \(x'\) are each \(F\)-related to \(y_0\).
Hence the formula is independent of both choices. Reversing \(R\)
constructs its inverse. Finally, \(xRy\) and \(uRv\) imply
\((xu)R(yv)\), so \(\theta\) is a homomorphism.

This proof is the small mechanism inside the comparison theorem:
\[
 \boxed{\text{a compatible relation}
 \ \longrightarrow\ \text{two ambiguity quotients}
 \ \longrightarrow\ \text{one canonical isomorphism}.}
\]
The abstract framework packages precisely the operations just used:
\[
\begin{array}{p{0.23\linewidth}|p{0.66\linewidth}}
 \textbf{Structure}&\textbf{Work it performs}\\ \hline
 Finite limits&Form pairs satisfying equations and the spaces of witnesses.\\
 Regular images&Form supports and compose relations.\\
 Barr exactness&Make equivalence relations into effective quotients.\\
 Mal'cev property&Supply the three-corner property for internal relations.\\
 Pointedness&Supply a zero object and zero classes for the later series argument.
\end{array}
\]
Pointedness is needed for that formulation of normal series, not for
the common-quotient construction itself. Groups and modules provide these
structures; the full exposition establishes the categorical relation
calculus before using it.

\pausehere{The conclusion is an isomorphism characterized by a commuting
diagram, not merely equality of the numbers of classes.}
\route{This page proves the group mechanism. The proof for every exact
Mal'cev category, with its image and quotient dependencies, is in
\mainref{1}{main §1} and \mainref{2}{main §2}.}

\clearpage
\section{Fill all four crossings}\label{c:rectangle}
Write \(A_0=B_0=0,\ A_1=A,\ B_1=B,\ A_2=B_2=G\).
For every pair of quotient steps form
\[
 P_{ij}=r_i s_j.
\]
Inputs come from the column side; outputs come from the row side.
Because partial equivalence relations are symmetric and idempotent,
\[
 P_{ij}^{\op}P_{ij}=s_jr_i s_j,\qquad
 P_{ij}P_{ij}^{\op}=r_i s_jr_i.
\]
The previous page supplies the canonical map from the input ambiguity
quotient to the output ambiguity quotient. Here every cell is explicit:
\[
\begin{array}{c|l|c|c}
 (i,j)&P_{ij}&\text{input quotient}&\text{output quotient}\\ \hline
 (1,1)&\{(0,0)\}&0/0&0/0\\
 (1,2)&\{(x,3x):x\in G\}&G/B&A/0\\
 (2,1)&\{(b,b),(b,b+3):b\in B\}&B/0&G/A\\
 (2,2)&G\times G&G/G&G/G .
\end{array}
\]
The two nontrivial isomorphisms are
\[
 \theta_{12}:G/B\xrightarrow{\sim}A,\quad x+B\longmapsto3x,
 \qquad
 \theta_{21}:B\xrightarrow{\sim}G/A,\quad b\longmapsto b+A.
\]
For \(P_{12}\), the intermediate point must lie in \(A\) and in the
\(B\)-coset of the input. It is \(0\) for even inputs and \(3\) for odd
inputs, exactly \(3x\). The formula descends because \(3B=0\).
Its inverse is \(a\mapsto a+B\): \(3a=a\) on \(A\), and
\(3x-x=2x\in B\).

For \(P_{22}\), any input \(x\) and output \(y\) have an intermediate
point \(z\) with \(z-x\in B\) and \(y-z\in A\), since \(A+B=G\).
Thus every pair is related. Both ambiguity quotients are trivial.
This cell has support \(G\), whereas the first cell has support \(0\).

Arrange the resulting factors by their original steps:
\[
\begin{array}{c|cc}
 &B/0&G/B\\ \hline
 A/0&0&C_2\\
 G/A&C_3&0 .
\end{array}
\]
The \(C_2\) cell compares the first \(A\)-step to the second \(B\)-step;
the \(C_3\) cell compares the other two steps.

\pausehere{A zero entry is the one-element group. It does not mean that
the relation is empty, or even that its support is small.}

There is still a logical obligation: the cells must assemble into
\emph{genuine refinements} of the original chains. A table of isomorphic
quotients by itself would not establish a refinement theorem.

\clearpage
\section{How cells become genuine refinements}\label{c:refinement}
In this abelian example the refined subgroups have familiar formulas:
\[
 U_{ij}=A_{i-1}+(A_i\cap B_j),\qquad
 V_{ji}=B_{j-1}+(B_j\cap A_i).
\]
For each fixed row or column, let the other index run from \(0\) to \(2\):
\[
\begin{array}{c|ccc}
 &j=0&j=1&j=2\\ \hline
 U_{1j}&0&0&A\\
 U_{2j}&A&G&G
\end{array}
\qquad
\begin{array}{c|ccc}
 &i=0&i=1&i=2\\ \hline
 V_{1i}&0&0&B\\
 V_{2i}&B&G&G .
\end{array}
\]
The output quotient in cell \((i,j)\) is
\(U_{ij}/U_{i,j-1}\); the input quotient is
\(V_{ji}/V_{j,i-1}\). Checking these against the previous table gives
the two refinements
\[
 0\subseteq0\subseteq A\subseteq G\subseteq G,\qquad
 0\subseteq0\subseteq B\subseteq G\subseteq G.
\]
Repeated terms contribute trivial factors and can later be deleted.

Why these intermediate subgroups? Let
\(q_i:A_i\to A_i/A_{i-1}\). Take the image of \(A_i\cap B_j\)
downstairs and pull it back:
\[
 q_i^{-1}\bigl(q_i(A_i\cap B_j)\bigr)
       =A_{i-1}+(A_i\cap B_j)=U_{ij}.
\]
Indeed \(q_i(x)=q_i(c)\) says exactly \(x-c\in A_{i-1}\).
The pullback includes \emph{every} representative of each chosen coset.
For example, in the second row the downstairs chain
\(0\subseteq G/A\subseteq G/A\) pulls back to
\(A\subseteq G\subseteq G\).

This explains \emph{saturation}: a relation downstairs must lift to all
representatives upstairs. For our \(R\subseteq B\times G\), changing an
output by \(3\) keeps the pair in \(R\); in formulas \(r_2R=R\).
Compare the smaller subgroup relation
\[
 T=\{(b,b):b\in B\}\subseteq B\times G.
\]
Under \(\id_B\times q_A\), both \(T\) and \(R\) have image
\(\{(b,b+A):b\in B\}\). But the inverse image of this graph is \(R\),
not \(T\): it includes \((b,b+3)\) too. Thus \(r_2T=R\ne T\).
Taking an image and pulling it back need not recover unsaturated data.

In the general relation proof, the row objects have the form
\(t=r_i s_jr_i\), and
\[
 r_i t r_i=t
\]
follows from \(r_i^2=r_i\). Saturation gives descent and pullback;
the boundary identities then make adjacent cells meet. The full proof
checks both requirements. Our two tables exhibit their result explicitly.

\route{For the categorical descent and boundary checks, see
\mainref{2}{main §2} and \mainref{3}{main §3}.}

\clearpage
\section{Why simplicity turns comparison into a matching}\label{c:simplicity}
Read the rectangle rowwise to get
\(0,C_2,C_3,0\); read it columnwise to get
\(0,C_3,C_2,0\). Delete the trivial factors. The matching is given by
actual maps:
\[
 A\xrightarrow{\sim}G/B,\quad a\longmapsto a+B,
 \qquad
 B\xrightarrow{\sim}G/A,\quad b\longmapsto b+A.
\]
No equality of lengths was needed to build these maps.

Here is the general final step for groups. A finite normal series in a
nontrivial simple group \(S\) has exactly one nontrivial factor.
Look at the first term equal to \(S\). Its preceding term is a proper
normal subgroup of \(S\), hence is trivial. Every earlier term is
therefore trivial, and every later term is \(S\). Notice that this argument
needs only normality in the \emph{next} term.

The general refinement theorem, proved in main §3, constructs a rectangle
of paired factors for two finite normal series. If the original series
are composition series, each row is a refinement of a simple quotient,
so exactly one nontrivial cell survives in that row. The same holds in
each column. Those surviving cells give a bijection between the original
factors, with an isomorphism for each matched pair. Equality of the numbers
of factors is a consequence of this bijection.

This booklet has computed that construction for \(C_6\) and proved its
common-quotient mechanism for groups. The full categorical refinement
theorem is supplied by the main exposition; a successful six-element
calculation is not a substitute for that theorem.

\textbf{Dimension now follows.}
For a finite basis \(v_1,\ldots,v_n\) of a vector space \(V\) over \(k\), put
\(F_i=\operatorname{span}(v_1,\ldots,v_i)\).
The map
\[
 k\longrightarrow F_i/F_{i-1},\qquad a\longmapsto av_i+F_{i-1}
\]
is onto by the definition of span and injective by linear independence.
The \(k\)-module \(k\) is simple: a submodule containing a nonzero \(a\)
contains \(a^{-1}a=1\), hence every scalar. A second finite basis therefore
gives another composition series, and Jordan--Hölder gives the same
number of steps. For \(V=0\), every basis is empty, directly by independence.

\textbf{What this example does not test.}
\(C_6\) is abelian, and the map \(A\times B\to G\), \((a,b)\mapsto a+b\),
is an isomorphism by §\ref{c:chains}. Thus all our subgroups are
ambient-normal and this extension splits. Exercise 6 tests adjacent
normality in \(A_4\). Exercise 7 shows that the same simple factors can
assemble into different groups.

\pausehere{This is a complete stopping point for the teaching route.
The higher chapters ask what the list of factors forgets, and what extra
structure can retain or recover that information.}

\clearpage
\section{Optional: the same comparison on successive pages}\label{c:spectral}
In the module setting, the two filtrations can instead be compared by
the spectral calculation developed in \mainref{7}{main §7}.
Here \(G\) is a \(\Z\)-module, and we use the two-term map
\(f:X=G\xrightarrow{\id}Y=G\).
Extend \(A_i,B_i\) by zero below \(0\) and by \(G\) above \(2\), and set
\[
 F_pX=A_{p-2},\qquad F_pY=B_p.
\]
The shift makes \(f\) filtration-preserving. On the initial associated
graded page, the only nonzero terms are
\[
\begin{array}{c|cccc}
 \text{filtration column}&1&2&3&4\\ \hline
 X&0&0&A&G/A\\
 Y&B&G/B&0&0 .
\end{array}
\]
The differential \(d_r\) goes from source column \(p\) to target column
\(p-r\). There is no same-column map, so \(d_0=0\).

\textbf{Page 1.}
The only nonzero differential is
\[
 d_{1,3}:A\longrightarrow G/B,\qquad a\longmapsto a+B.
\]
It is the \(C_2\) comparison already proved to be an isomorphism.
Taking its kernel and cokernel removes both of these entries.

\textbf{Page 2.}
The remaining source is \(G/A\) in column \(4\), and the remaining
target is \(B\) in column \(1\). A page-2 differential moves two columns,
so \(d_2=0\). One zero differential does not imply stabilization.

\textbf{Page 3.}
The representative subgroup for the surviving source is now \(B\):
to land three filtration steps lower, an element of \(G\) must lie in
\(B\). Its identification with the preceding source term is precisely
\[
 B\xrightarrow{\sim}G/A,\qquad b\longmapsto b+A.
\]
On these representatives \(f\) is the identity \(B\to B\), so
\(d_{3,4}\) is an isomorphism. If we keep the old \(G/A\) coordinate,
the same map is \(x+A\mapsto4x\). These are two coordinate descriptions
of one differential, not two different maps. Page 4 is zero.

The calculation pairs the \(C_2\)'s first and the \(C_3\)'s later.
It uses the same quotient maps as the rectangle; their page locations
depend on the chosen filtrations and shift. The filtration machinery
and its convergence are proved in the main exposition, not assumed from
Jordan--Hölder. Here we have computed the pages explicitly.

\route{Read this after the concrete rectangle. It is an abelian
\(\Z\)-module calculation; it does not assert that ordinary abelian
spectral sequences apply directly to arbitrary groups.}
